% Corrigé intégré automatiquement pour la banque Exercices.
% Source : A_Integrer/DDS_04/071_PFS/071_PFS.tex
\begin{corrigebox}[Corrigé]
\CorrigeQuestion{1}
En considérant le poids de la bobine de papier seul, on a directement 
$P_B (R)= \rho_v L \pi \left(R^2 - R_M^2 \right) g$
\CorrigeQuestion{2}
$P_B = \SI{101764}{N}$.
\CorrigeQuestion{3}
On isole la bobine.

BAME : poids de la bobine $\vect{P_B}$, efforts presseurs des rouleaux en $A$ et en $B$, tension $\vect{T}$ de la feuille en $A$, effort du rouleau presseur.

TRS : $\vect{F(t)} + \vect{R_A} + \vect{R_B} + \vect{P_B} + \vect{T} = \vect{0}$
\CorrigeQuestion{4}
Les efforts des rouleaux porteurs se décomposent en un effort tangentiel et un effort normal, orienté vers $C_B$.

\begin{center}
\includegraphics[width=.35\linewidth]{988_03}%
\end{center}
\CorrigeQuestion{5}
On isole le rouleau porteur arrière.

BAME : action de la bobine fille en $B$, action de la pesanteur sur la bobine fille en $C_2$, couple moteur en $C_2$. 

Si on néglige l'inertie de la bobine arrière, on peut appliquer le théorème du moment statique en $C_2$ en projection sur $\vect{z}$ :
\begin{itemize}
\item $\vectm{C_2}{\text{fille}}{\text{porteur ar}} = \vect{C_2 B} \wedge \vect{R_B} = \pm r T_B$; 
\item $\vectm{C_2}{\text{pes}}{\text{porteur ar}} = \vect{0}$; 
\item $C_m$.
\end{itemize}
Au final $C_m = \pm r T_B$.
\CorrigeQuestion{6}
On projette l'expression de la question 4 sur $\vect{x}$ et on a 

$R_A \sin \left(\theta+\varphi\right) - R_B\sin \left(\theta+\varphi\right) + T \cos\theta = 0$

soit
$T =\left(R_B  -R_A\right) \dfrac{\sin \left(\theta+\varphi\right)}{\cos\theta } $
\CorrigeQuestion{7}
On projette l'expression de la question 4 sur $\vect{y}$ et on a 

\noindent $R_A \cos \left(\theta+\varphi\right) - R_B\cos \left(\theta+\varphi\right) - T \sin\theta -F(t)-P= 0$ 

soit

$F(t)=R_A \cos \left(\theta+\varphi\right) - R_B\cos \left(\theta+\varphi\right) - T \sin\theta -P$.

En utilisant le résultat de la question précédente, on a : 

$F(t)=R_A \cos \left(\theta+\varphi\right) - R_B\cos \left(\theta+\varphi\right) - \left(R_B  -R_A\right) \dfrac{\sin \left(\theta+\varphi\right)}{\cos\theta } \sin\theta -P$.

et

$F(t)=R_A \cos \left(\theta+\varphi\right) - R_B\cos \left(\theta+\varphi\right) - \left(R_B  -R_A\right)\sin \left(\theta+\varphi\right)  \tan \theta -P$.
\end{corrigebox}
